\documentclass[a4paper]{article}
\usepackage{amsmath}
\usepackage{amssymb}
\usepackage{mathtools}
\usepackage[margin=0.5in]{geometry}

\begin{document}

\section{Questão}

Seja

\begin{equation}
    \mathbf{T} = \begin{bmatrix}
        \mathbf{Q} & \mathbf{s} \\
        \mathbf{0}_{1\times 3} & 1
    \end{bmatrix}
\end{equation}

uma matriz de transformação homogênea, onde $\mathbf{Q}$ é uma matriz de rotação $3 \times 3$ e $\mathbf{s}$ é um vetor de translação $3 \times 1$.

Há uma matriz $4 \times 4$, $\mathbf{M}$, e um escalar, $\alpha$, tais que $\alpha \in \mathbb{R}$, $\mathbf{M} \neq \mathbf{T}$, e

\begin{equation}
    \mathbf{T}\cdot\begin{bmatrix}\mathbf{p} \\ 1\end{bmatrix} = \mathbf{M}\cdot\begin{bmatrix}\mathbf{p} \\ \alpha\end{bmatrix}
\end{equation}

se verifica para todo vetor $\mathbf{p} \in \mathbb{R}^3$?

\section{Investigação}

Seja $\mathbf{M} = \begin{bmatrix}
    \mathbf{A} & \mathbf{b} \\
    \mathbf{c}^T & d
\end{bmatrix}$, onde $\mathbf{A}$ é uma matriz $3 \times 3$, $\mathbf{b}$ é um vetor $3 \times 1$, $\mathbf{c}$ é um vetor $3 \times 1$, e $d$ é um escalar.

Se a condição dada é verdadeira,

\begin{align}
    \begin{bmatrix}
        \mathbf{Q} & \mathbf{s} \\
        \mathbf{0}_{1\times 3} & 1
    \end{bmatrix}\cdot\begin{bmatrix}\mathbf{p} \\ 1\end{bmatrix} &=
    \begin{bmatrix}
        \mathbf{A} & \mathbf{b} \\
        \mathbf{c}^T & d
    \end{bmatrix}\cdot\begin{bmatrix}\mathbf{p} \\ \alpha\end{bmatrix} \\
    \begin{bmatrix}
        \mathbf{Q}\mathbf{p} + \mathbf{s} \\
        1
    \end{bmatrix} &=
    \begin{bmatrix}
        \mathbf{A}\mathbf{p} + \alpha\mathbf{b} \\
        \mathbf{c}^T\mathbf{p} + \alpha d
    \end{bmatrix}
\end{align}

Para que isso seja verdade para todo $\mathbf{p}$,

\begin{align}
    \mathbf{Q}\mathbf{p} + \mathbf{s} &=
    \mathbf{A}\mathbf{p} + \alpha\mathbf{b} \label{eq:vetor_0} \\
    1 &=
    \mathbf{c}^T\mathbf{p} + \alpha d \label{eq:vetor_1}
\end{align}

De (\ref{eq:vetor_1}), temos que $\mathbf{c}^T\mathbf{p} = 1 - \alpha d$. Para que isso seja verdade para todo $\mathbf{p}$, $\mathbf{c}$ deve ser o vetor nulo e $d$ deve ser igual a $1/\alpha$.

De (\ref{eq:vetor_0}), temos que $\mathbf{Q}\mathbf{p} + \mathbf{s} = \mathbf{A}\mathbf{p} + \alpha\mathbf{b}$, que requer $\mathbf{Q} = \mathbf{A}$ e $\mathbf{b} = \mathbf{s}/\alpha$.

Portanto, para $\alpha \in \mathbb{R} \backslash \{0, 1\}$, $\mathbf{c} = \mathbf{0}$, $d = 1/\alpha$, $\mathbf{A} = \mathbf{Q}$, e $\mathbf{b} = \mathbf{s}/\alpha$, temos que

\begin{equation}
    \begin{cases}
        \mathbf{M} \neq \mathbf{T} \\
        \mathbf{T}\cdot\begin{bmatrix}\mathbf{p} \\ 1\end{bmatrix} =
        \mathbf{M}\cdot\begin{bmatrix}\mathbf{p} \\ \alpha\end{bmatrix}
    \end{cases}
\end{equation}

São $\mathbf{M}$ e $\mathbf{T}$ linearmente independentes?

Se forem dependentes, $\mathbf{M} = k\mathbf{T}$ para algum $k \in \mathbb{R}$.

\begin{equation}
    \begin{bmatrix}
        \mathbf{A} & \mathbf{b} \\
        \mathbf{0}_{1\times 3} & 1/\alpha
    \end{bmatrix} =
    \begin{bmatrix}
        k\mathbf{Q} & k\mathbf{s} \\
        \mathbf{0}_{1\times 3} & k
    \end{bmatrix}
\end{equation}

Dado que, de (\ref{eq:vetor_0}), $\mathbf{A} = \mathbf{Q}$ e $\mathbf{b} = \mathbf{s}/\alpha$, temos que $k = 1$ e $k = 1/\alpha$. Isso é uma contradição, pois $\alpha \neq 1$. Portanto, $\mathbf{M}$ e $\mathbf{T}$ são linearmente independentes.

\section{Conclusão}

Sim, existe uma matriz $4 \times 4$, $\mathbf{M}$, e um escalar, $\alpha$, tais que $\alpha \in \mathbb{R}$, $\mathbf{M} \neq \mathbf{T}$, e

\begin{equation}
    \mathbf{T}\cdot\begin{bmatrix}\mathbf{p} \\ 1\end{bmatrix} = \mathbf{M}\cdot\begin{bmatrix}\mathbf{p} \\ \alpha\end{bmatrix}
\end{equation}

se verifica para todo vetor $\mathbf{p} \in \mathbb{R}^3$.

Além disso, $\mathbf{M}$ e $\mathbf{T}$ são linearmente independentes.

Qual a aplicação disso? Não sei.

\end{document}
